\documentclass{article}
\usepackage{amssymb}
\usepackage{amsmath}
\usepackage{forest}

\title{TRADECRAFT BLUEPRINT v0.1}
\author{Samim C. Sulog}
\date{\today}

\begin{document}
\maketitle
This is a model paper that models the tradecraft system in the game theory language, utilising the primitives of the game (the core components) and structuring based on game design principles.

\section{Environment}
The primitives of the game, or the core components of what builds the tradecraft system itself, from \textbf{players}, \textbf{strategy space}, \textbf{payoff functions}, \textbf{operator space}, and \textbf{information sets or history}. All of this will be explained in a game design lens as well.

\subsection{Players}
The player component is also called in game design terms, the \textbf{user}. The mathematical expression of the variable is:
\[
N = \{\text{Agent}, \text{Target}\}
\]

\subsection{Operator Space}
The operator space is also called the \textbf{game transformation} in game design terms. The mathematical expression is broken down into sub-components that build the core of tradecraft:
\[
O = \{O_{\text{Recon}}, O_{\text{Manip}}, O_{\text{Protect}}, O_{\text{Offense}}\}
\]

\subsubsection{Information Reconaissance}
\[
O_{\text{Recon}}(I_{\text{Agent}}) = \{h_{\text{true}}\}
\]
\begin{small}
Information Reconaissance is part of intelligence work, notably, the process of the Intelligence Cycle is what exactly builds and defines the component of Reconaissance, the reconaissance operator collapses the Info Set of a player to absolute truth at the current moment, where the absolute truth is defined as $h_{\text{true}}$. 
\end{small}

\subsubsection{Information Protection}
\begin{align*}
O_{\text{Protect}}(I_{\text{Target}}) = I_{\text{enemy}} \cup \{\text{Decoy Node}\} \\
O_{\text{Protect}}(u_{\text{Target}}(s)) = u_{\text{Target}}(s) - \Delta \text{Cost}
\end{align*}
\begin{small}
The process of protecting one's information is in line with intelligence work, especially with the Operation Security Cycle, in game theory terms, this is more in line with expanding the enemy's Info Set with fake nodes, or increasing the cost necessary to the point the attack itself is unnecessary.
\end{small}

\subsubsection{Information Manipulation}
\[
O_{\text{Manip}}(u_{\text{Target}}) =u^{'}_{\text{Target}}
\]
\begin{small}
The process of manipulating information, is also called \textbf{forgery}. It is also used in intelligence work to deceive the enemy's spy. Which is also intersecting the Protection pillar, especially the decoy node, but in this operator, it is also injecting a false payoff to make the enemy think a losing move is a winning move.
\end{small}

\subsubsection{Information Offense}
\[
O_{\text{Offense}}(S_{\text{Target}}) = \frac{S_{\text{Target}}}{\text{Strategy to Remove}}
\]
\begin{small}
The use of information as an offense or attack is also part of intelligence work, the offense vary from character assassination, to misinformation, to radicalization campaign, and to hacking. In game theory terms, this is more or so removing the target's capability to do certain strategies or moves.
\end{small}

\subsection{Strategy Space}
The strategy space is also called the \textbf{system} in game design terms. The mathematical expression is broken down into sub-components that build the system of tradecraft:
\begin{align*}
S_{\text{Agent}} &= \{O\} \\
S_{\text{Target}} &= \{O\} \\
S &= S_{\text{Agent}} \cdot S_{\text{Target}} \\
\text{System}_{\text{Agent}} &= \{S_{\text{Agent}}, u_{\text{Agent}}\} \\
\text{System}_{\text{Target}} &= \{S_{\text{Target}}, u_{\text{Target}}\} \\
\text{System} &= \{S, \{u\}\} = \text{System}_{\text{Agent}} \cdot \text{System}_{\text{Target}}
\end{align*}
Each mathematical expression builds up from one another to define the System variable. In game theory, the term \textbf{system} is also synonymous with the term \textbf{game} due to how equivalent they are mathematically, but for rigor and clarity, system is what builds the game; therefore, defining system inside the game set clarifies things rigorously.

\subsection{Payoff Functions}
The payoff functions are also called the \textbf{mechanics} in game design terms. The mathematical expression is broken down into sub-components to build the mechanics of tradecraft:
\begin{align*}
\text{Advantage}_{\text{Agent}} = |I_{\text{Agent}}| - |I_{\text{Spy}}| \\
P_{\text{success}} = \min(1, \max(0, 0.5 + (Advantage_{\text{Spy}} \times 0.1))) \\
u_{\text{Spy}} = (P_{\text{success}} \times \text{Win Payoff}) + ((1 - P_{\text{success}}) \times \text{Loss Payoff}) \\
u_{\text{Target}} = -u_{\text{Spy}}
\end{align*}
The payoff function is calibrated with the information advantage, as the more advantage you have with the information, the more certain your strategies are.

\subsection{Information Sets}
The information set in game theory and game design, are the datas that is either available to the user, or some user, or not at all; in game design terms, this is called the \textbf{Level of Knowledge}. The meaning of information may be the situation, the description of an item, or a move that was just done. Game theory wise, this is denoted by nodes in a game tree:
\begin{align*}
  I &= \{ \text{Node}_1, \text{Node}_2 \}
\end{align*}
In intelligence works, the information themselves are imperfect information, hence the process of reconaissance. The game itself is imperfect information due to also the process of protection, and manipulation. Certain information needs to be checked before they are assumed true via \textbf{Admiralty Code} and \textbf{CRAP} under \textbf{Process of the Intelligence Cycle}.

\begin{center}
\resizebox{\linewidth}{!}{% Scales the tree down only if needed, or remove this line if s sep=8mm fits on its own
\begin{forest}
  % 1. Style for player/nature dots
  dot/.style={circle, fill=black, inner sep=1.5pt, outer sep=0pt},
  % 2. Shortcut for action labels on branches: act={position}{text}
  act/.style 2 args={edge label={node[midway, #1, font=\small]{#2}}},
  for tree={
    l sep=16mm, % Vertical branch length
    s sep=8mm,  % Reduced horizontal spread to prevent margin overflow
    font=\small % Slightly smaller text for payoff nodes
  }
  % --- TREE STRUCTURE ---
  [, dot, label=above:{Target}
    [, dot, label=left:{Node}, act={left}{Leave Trace}
      [, dot, name=agent1, label=above left:{Agent}, act={above left}{Real}
        [{($u_{\text{Target}}$, $u_{\text{Agent}}$)}, act={left}{Steal}] 
        [{($u_{\text{Target}}$, $u_{\text{Agent}}$)}, act={right}{Recon}]
      ]
      [, dot, name=agent2, label=above right:{Agent}, act={above right}{Fake}
        [{($u_{\text{Target}}$, $u_{\text{Agent}}$)}, act={left}{Steal}]  
        [{($u_{\text{Target}}$, $u_{\text{Agent}}$)}, act={right}{Recon}]
      ]
    ]
  ]
  % --- IMPERFECT INFORMATION SET ---
  \draw[dashed, thick] (agent1) -- (agent2);
\end{forest}%
}
\end{center}

\section{The Complete Game}
\begin{align*}
  N &= \{\text{Agent}, \text{Target}\} \\
  S &= S_{\text{Agent}} \times S_{\text{Target}} \\
  u &= \big(u_{\text{Agent}},\, S_{\text{Target}}\big) \\
  \text{System} &= \big(S,\, u\big) \\
  I_{n}\subseteq H &= \text{Player's Information Set} \\
  H &= \{h_1, h_2, \dots, h_i\} \\
  G(L) &= \big(N,\, \text{System},\, I,\, H\big)
\end{align*}
Law is a game of rules, but Tradecraft is a game of Information. In game theory terms, Tradecraft is a \textbf{zero-sum game of imperfect and incomplete information.}


\end{document}
